\documentclass[reqno, answers]{exam}

\usepackage{amssymb, amsmath, amsthm, stmaryrd}
\usepackage{fullpage, parskip}
\usepackage{enumitem}
\usepackage{tikz}
\usepackage{ulem}
\usepackage{hyperref}

\title{Introduction to Computational Math Spring 2026 \\ Homework 05}
\date{Due: Friday, Feb 27, 2026, 11:59 PM}
\author{}

\begin{document}
\maketitle
\thispagestyle{empty}

Textbook = Driscoll and Braun (\url{https://fncbook.com/})

There are more problems in this homework than the previous one, but they are shorter and more conceptual.


\begin{enumerate}

%% Exercise 3.2.1
\item Textbook Exercise 3.2.1: Work out the least-squares solution when
\[
A = \begin{bmatrix} 2 & -1 \\ 0 & 1 \\ -2 & 2 \end{bmatrix}, \quad b = \begin{bmatrix} 1 \\ -5 \\ 6 \end{bmatrix}.
\]

\begin{solution}
The least-squares solution is given by $x = (A^T A)^{-1} A^T b$.

First, compute $A^T A$:
\[
A^T A = \begin{bmatrix} 2 & 0 & -2 \\ -1 & 1 & 2 \end{bmatrix} \begin{bmatrix} 2 & -1 \\ 0 & 1 \\ -2 & 2 \end{bmatrix} = \begin{bmatrix} 4+0+4 & -2+0-4 \\ -2+0-4 & 1+1+4 \end{bmatrix} = \begin{bmatrix} 8 & -6 \\ -6 & 6 \end{bmatrix}
\]

Next, compute $A^T b$:
\[
A^T b = \begin{bmatrix} 2 & 0 & -2 \\ -1 & 1 & 2 \end{bmatrix} \begin{bmatrix} 1 \\ -5 \\ 6 \end{bmatrix} = \begin{bmatrix} 2 + 0 - 12 \\ -1 - 5 + 12 \end{bmatrix} = \begin{bmatrix} -10 \\ 6 \end{bmatrix}
\]

Now find $(A^T A)^{-1}$. The determinant is $\det(A^T A) = 8 \cdot 6 - (-6)(-6) = 48 - 36 = 12$.
\[
(A^T A)^{-1} = \frac{1}{12} \begin{bmatrix} 6 & 6 \\ 6 & 8 \end{bmatrix} = \begin{bmatrix} 1/2 & 1/2 \\ 1/2 & 2/3 \end{bmatrix}
\]

Finally, compute the solution:
\[
x = (A^T A)^{-1} A^T b = \begin{bmatrix} 1/2 & 1/2 \\ 1/2 & 2/3 \end{bmatrix} \begin{bmatrix} -10 \\ 6 \end{bmatrix} = \begin{bmatrix} -5 + 3 \\ -5 + 4 \end{bmatrix} = \begin{bmatrix} -2 \\ -1 \end{bmatrix}
\]
\end{solution}

%% Exercise 3.2.2
\item Textbook Exercise 3.2.2: Use equation (3.2.4) to find the pseudoinverse $A^+$ of the matrix $A = \begin{bmatrix} 1 & -2 & 3 \end{bmatrix}^T$.

\begin{solution}
The pseudoinverse is defined as $A^+ = (A^T A)^{-1} A^T$ for a matrix with full column rank.

Here $A = \begin{bmatrix} 1 \\ -2 \\ 3 \end{bmatrix}$ is a $3 \times 1$ matrix.

Compute $A^T A$:
\[
A^T A = \begin{bmatrix} 1 & -2 & 3 \end{bmatrix} \begin{bmatrix} 1 \\ -2 \\ 3 \end{bmatrix} = 1 + 4 + 9 = 14
\]

Thus $(A^T A)^{-1} = \frac{1}{14}$.

The pseudoinverse is:
\[
A^+ = (A^T A)^{-1} A^T = \frac{1}{14} \begin{bmatrix} 1 & -2 & 3 \end{bmatrix} = \begin{bmatrix} \frac{1}{14} & -\frac{2}{14} & \frac{3}{14} \end{bmatrix} = \begin{bmatrix} \frac{1}{14} & -\frac{1}{7} & \frac{3}{14} \end{bmatrix}
\]
\end{solution}

%% Exercise 3.2.3
\item Textbook Exercise 3.2.3: Prove the first statement of Theorem 3.2.2: $A^T A$ is symmetric for any $m \times n$ matrix $A$ with $m \geq n$.

\begin{solution}
To show that $A^T A$ is symmetric, we need to prove that $(A^T A)^T = A^T A$.

Using the property $(XY)^T = Y^T X^T$:
\[
(A^T A)^T = A^T (A^T)^T = A^T A
\]

Since $(A^T A)^T = A^T A$, the matrix $A^T A$ is symmetric.
\end{solution}

%% Exercise 3.2.5a
\item Textbook Exercise 3.2.5a: Show that for any $m \times n$ matrix $A$ with $m > n$ for which $A^T A$ is nonsingular, $A^+ A$ is the $n \times n$ identity.

\begin{solution}
The pseudoinverse is defined as $A^+ = (A^T A)^{-1} A^T$.

Computing $A^+ A$:
\[
A^+ A = (A^T A)^{-1} A^T A
\]

Since $A^T A$ is nonsingular (invertible), we have:
\[
A^+ A = (A^T A)^{-1} (A^T A) = I_n
\]

where $I_n$ is the $n \times n$ identity matrix.
\end{solution}

%% Exercise 3.4.1
\item Textbook Exercise 3.4.1: Find a Householder reflector $P$ such that
\[
P \begin{bmatrix} 2 \\ 9 \\ -6 \end{bmatrix} = \begin{bmatrix} 11 \\ 0 \\ 0 \end{bmatrix}.
\]

\begin{solution}
A Householder reflector has the form $P = I - 2\frac{vv^T}{v^T v}$, where $v = z - \|z\|_2 e_1$ (or $v = z + \|z\|_2 e_1$).

Let $z = \begin{bmatrix} 2 \\ 9 \\ -6 \end{bmatrix}$.

Compute $\|z\|_2 = \sqrt{4 + 81 + 36} = \sqrt{121} = 11$.

Using $v = z - \|z\|_2 e_1$:
\[
v = \begin{bmatrix} 2 \\ 9 \\ -6 \end{bmatrix} - 11 \begin{bmatrix} 1 \\ 0 \\ 0 \end{bmatrix} = \begin{bmatrix} -9 \\ 9 \\ -6 \end{bmatrix}
\]

Compute $v^T v = 81 + 81 + 36 = 198$.

The Householder reflector is:
\[
P = I - 2\frac{vv^T}{v^T v} = I - \frac{2}{198} \begin{bmatrix} -9 \\ 9 \\ -6 \end{bmatrix} \begin{bmatrix} -9 & 9 & -6 \end{bmatrix}
\]

\[
P = I - \frac{1}{99} \begin{bmatrix} 81 & -81 & 54 \\ -81 & 81 & -54 \\ 54 & -54 & 36 \end{bmatrix} = \begin{bmatrix} 1 - \frac{81}{99} & \frac{81}{99} & -\frac{54}{99} \\ \frac{81}{99} & 1 - \frac{81}{99} & \frac{54}{99} \\ -\frac{54}{99} & \frac{54}{99} & 1 - \frac{36}{99} \end{bmatrix}
\]

\[
P = \begin{bmatrix} \frac{18}{99} & \frac{81}{99} & -\frac{54}{99} \\ \frac{81}{99} & \frac{18}{99} & \frac{54}{99} \\ -\frac{54}{99} & \frac{54}{99} & \frac{63}{99} \end{bmatrix} = \frac{1}{11} \begin{bmatrix} 2 & 9 & -6 \\ 9 & 2 & 6 \\ -6 & 6 & 7 \end{bmatrix}
\]
\end{solution}

%% Exercise 3.4.2
\item Textbook Exercise 3.4.2: Prove the unfinished items in Theorem 3.4.1, namely, that a Householder reflector $P$ is symmetric and orthogonal.

To show that a square matrix $Q$ is orthogonal, it is enough to show that $Q^T Q = I$.

\begin{solution}
The Householder reflector is defined as:
\[
P = I - 2\frac{vv^T}{v^T v}
\]

\textbf{Proof that $P$ is symmetric:}

We need to show $P^T = P$.
\[
P^T = \left(I - 2\frac{vv^T}{v^T v}\right)^T = I^T - 2\frac{(vv^T)^T}{v^T v} = I - 2\frac{vv^T}{v^T v} = P
\]
since $(vv^T)^T = vv^T$. Thus $P$ is symmetric.

\textbf{Proof that $P$ is orthogonal:}

Since $P$ is symmetric, $P^T = P$. We need to show $P^T P = P^2 = I$.
\[
P^2 = \left(I - 2\frac{vv^T}{v^T v}\right)\left(I - 2\frac{vv^T}{v^T v}\right)
\]

Expanding:
\[
P^2 = I - 4\frac{vv^T}{v^T v} + 4\frac{vv^T vv^T}{(v^T v)^2}
\]

Note that $vv^T v = v(v^T v)$, so:
\[
\frac{vv^T vv^T}{(v^T v)^2} = \frac{v(v^T v)v^T}{(v^T v)^2} = \frac{vv^T}{v^T v}
\]

Therefore:
\[
P^2 = I - 4\frac{vv^T}{v^T v} + 4\frac{vv^T}{v^T v} = I
\]

Thus $P^T P = P^2 = I$, confirming $P$ is orthogonal.
\end{solution}

%% Exercise 3.4.3
\item Textbook Exercise 3.4.3: Let $P$ be a Householder reflector as in (3.4.1).
\begin{enumerate}
\item[(a)] Find a vector $u$ such that $Pu = -u$.
\item[(b)] What algebraic condition is necessary and sufficient for a vector $x$ to satisfy $Px = x$? In $n$ dimensions, how many such linearly independent vectors are there?
\end{enumerate}

\begin{solution}
Recall $P = I - 2\frac{vv^T}{v^T v}$.

\textbf{Part (a):} We want to find $u$ such that $Pu = -u$.
\[
Pu = \left(I - 2\frac{vv^T}{v^T v}\right)u = u - 2\frac{v(v^T u)}{v^T v} = -u
\]

This requires:
\[
2u = 2\frac{v(v^T u)}{v^T v}
\]

Taking $u = v$ works:
\[
Pv = v - 2\frac{v(v^T v)}{v^T v} = v - 2v = -v
\]

So $u = v$ (or any scalar multiple of $v$) satisfies $Pu = -u$.

\textbf{Part (b):} We want $Px = x$, which means:
\[
Px = x - 2\frac{v(v^T x)}{v^T v} = x
\]

This requires $\frac{v(v^T x)}{v^T v} = 0$, which means $v^T x = 0$.

The algebraic condition is: $v^T x = 0$, i.e., $x$ must be orthogonal to $v$.

In $n$ dimensions, the orthogonal complement of $v$ has dimension $n-1$. Therefore, there are $n-1$ linearly independent vectors satisfying $Px = x$.
\end{solution}

%% Exercise 3.4.7a
\item Textbook Exercise 3.4.7a: Given a 2-vector $[\alpha, \beta]^T$, a Givens rotation defines an angle $\theta$ such that
\[
\begin{bmatrix} \cos\theta & \sin\theta \\ -\sin\theta & \cos\theta \end{bmatrix} \begin{bmatrix} \alpha \\ \beta \end{bmatrix} = \begin{bmatrix} \sqrt{\alpha^2 + \beta^2} \\ 0 \end{bmatrix}.
\]
Given $\alpha$ and $\beta$, show how to compute $\cos\theta$ and $\sin\theta$ without evaluating any trig functions.

This question is asking you to find $\theta$ in terms of $\alpha$ and $\beta$. Your answer will involve a trigonometric function.

\begin{solution}
Let $r = \sqrt{\alpha^2 + \beta^2}$.

From the matrix equation:
\[
\begin{bmatrix} \cos\theta & \sin\theta \\ -\sin\theta & \cos\theta \end{bmatrix} \begin{bmatrix} \alpha \\ \beta \end{bmatrix} = \begin{bmatrix} \alpha\cos\theta + \beta\sin\theta \\ -\alpha\sin\theta + \beta\cos\theta \end{bmatrix} = \begin{bmatrix} r \\ 0 \end{bmatrix}
\]

From the second equation:
\[
-\alpha\sin\theta + \beta\cos\theta = 0 \quad \Rightarrow \quad \beta\cos\theta = \alpha\sin\theta
\]

From the first equation:
\[
\alpha\cos\theta + \beta\sin\theta = r
\]

From $\beta\cos\theta = \alpha\sin\theta$, if we also use $\cos^2\theta + \sin^2\theta = 1$, we can solve directly.

Squaring and adding the relations derived from the geometry, we get:
\[
\cos\theta = \frac{\alpha}{\sqrt{\alpha^2 + \beta^2}} = \frac{\alpha}{r}
\]
\[
\sin\theta = \frac{\beta}{\sqrt{\alpha^2 + \beta^2}} = \frac{\beta}{r}
\]

These formulas allow us to compute $\cos\theta$ and $\sin\theta$ directly from $\alpha$ and $\beta$ without evaluating any trigonometric functions---only arithmetic operations and a square root are needed.
\end{solution}

\end{enumerate}

No Jupyter notebook this week.



\section*{Quiz Information}

The quiz will be held on {\bf Tuesday, March 03, 2026} during the discussion section.

Quiz questions will be modifications of the homework problems and material discussed in class.

\textbf{Topics covered:} 

Chapters 3.2 and 3.4. 

You will need memorize 
\begin{itemize}
    \item The normal equations for least squares problems,
    \item The formula for the Householder reflector.
\end{itemize}

You should expect both computational and conceptual questions on the quiz.

There won't be any questions about using the normal equations in terms of $QR$ factorization, but you should know how to use the normal equations in terms of the $A^T A$ matrix.



\end{document}